For embedded and edge systems, logging must balance observability with overhead (storage, CPU, and concurrent initialization). This thesis implements Optbinlog, a schema-shared binary logging approach for embedded deployment. Event metadata are parsed once, materialized in shared memory, and accessed via offsets, so runtime hot-path cost is reduced and cross-process mapping remains stable.

The work first surveys representative domestic and international logging systems (Linux ftrace, LTTng, Zephyr logging, systemd-journald, NanoLog, RT-Thread ulog, OpenHarmony HiLog Lite) and analyzes their optimization code paths. Then a theory-first model is built to derive how time/space/throughput scale with log count \texttt{N} and device count \texttt{D}. Finally, measured results are compared against the derived trends.

The retained dataset follows a stage-based workflow: engineering-approximate matrix rerun, strict-aligned suite (single and local multi-device scans at \texttt{D=5,10,20,50}), hardware-CRC32C real multi-node rerun (\texttt{5/10/15/20} nodes), and local/node-level initialization competition. Results show that \texttt{binary} consistently outperforms \texttt{text} across the included semantic groups. Under strict semantic alignment, Optbinlog remains stably ahead of \texttt{ulog/hilog/syslog}, while \texttt{nanolog/zephyr} still show persistent node-level penalties in real multi-node runs. In addition, \texttt{ftrace} in real multi-node runs shows a split pattern (space-positive but time/throughput-negative). The added space-scan experiment shows early crossover for \texttt{ulog/hilog/syslog} at around \texttt{N≈10}, while no crossover is observed for \texttt{nanolog/zephyr} up to \texttt{N=100000}. Initialization-race results satisfy the single-creator property consistently. Overall, theoretical trends and empirical observations are aligned.
